\subsubsection{Teacher--Subject সম্পর্ক ও সংযোগ Table}
% Parent -> junction <- parent, using the topic's unchanged source records.
\begin{flushleft}
\begin{tikzpicture}[every node/.style={font=\footnotesize,inner sep=0pt,align=left}]
\node (teacher) {\begin{tabular}{|l|l|}\hline
\multicolumn{2}{|l|}{\textbf{Teacher}}\\\hline
\chSixHeader TeacherID & Name\\\hline
T1001 & Hasan\\\hline T1002 & Nila\\\hline
T1003 & Rafi\\\hline T1004 & Mitu\\\hline
\end{tabular}};
\node[anchor=north west] (teaching) at ([xshift=11mm]teacher.north east)
{\begin{tabular}{|l|l|l|}\hline
\multicolumn{3}{|l|}{\textbf{Teaching}}\\\hline
\chSixHeader TeacherID & SubjectCode & Day\\\hline
T1001 & 101 & Sat\\\hline T1001 & 102 & Sun\\\hline
T1002 & 101 & Mon\\\hline T1003 & 103 & Tue\\\hline
T1003 & 104 & Wed\\\hline
\end{tabular}};
\node[anchor=north west] (subject) at ([xshift=11mm]teaching.north east)
{\begin{tabular}{|l|l|}\hline
\multicolumn{2}{|l|}{\textbf{Subject}}\\\hline
\chSixHeader SubjectCode & Subject\\\hline
101 & Bangla\\\hline 102 & English\\\hline
103 & ICT\\\hline 104 & Physics\\\hline
\end{tabular}};
\draw[bookarrow] (teacher.east)--node[pos=.2,above=2pt]{1}node[pos=.8,above=2pt]{M}(teaching.west);
\draw[bookarrow] (subject.west)--node[pos=.2,above=2pt]{1}node[pos=.8,above=2pt]{M}(teaching.east);
\node[anchor=north west,text width=29mm] at ([yshift=-2mm]teacher.south west) {TeacherID = PK};
\node[anchor=north west,text width=45mm] at ([yshift=-2mm]teaching.south west) {TeacherID ও SubjectCode = FK};
\node[anchor=north west,text width=32mm] at ([yshift=-2mm]subject.south west) {SubjectCode = PK};
\end{tikzpicture}
\captionof{table}{দুই Parent থেকে Teaching Junction Table-এ 1:M সংযোগ}
\end{flushleft}

\begin{multicols}{2}
\textbf{Teacher থেকে পড়ি।} Hasan-এর ID T1001। Teaching-এ এই ID-র
সঙ্গে 101 ও 102 রয়েছে; Subject থেকে নাম মিলিয়ে পাই Bangla ও English।
অর্থাৎ একজন শিক্ষক একাধিক বিষয় পড়ান।

\columnbreak
\textbf{Subject থেকে পড়ি।} Bangla-এর Code 101। Teaching-এ 101-এর
সঙ্গে T1001 ও T1002 রয়েছে। Teacher থেকে পাই Hasan ও Nila। অর্থাৎ
একটি বিষয় একাধিক শিক্ষক পড়ান।
\end{multicols}

Teaching-এর প্রতিটি সারি একজন শিক্ষক ও একটি বিষয়ের সংযোগ এবং তার দিন
রাখে। TeacherID ও SubjectCode আলাদা Foreign Key; একটি মানের পুনরাবৃত্তি
থাকলেও তা একই রেকর্ডের অনিচ্ছাকৃত অনুলিপি নয়। এই নমুনায় একই
শিক্ষক--বিষয় জোড়ার জন্য একটি দিন রাখা হয়েছে। একই জোড়া একবার রাখার
নিয়ম থাকলে জোড়াটিকে Composite Primary Key করা যায়।

Teacher থেকে Teaching এবং Subject থেকে Teaching---দুটি সম্পর্কই 1:M।
এই দুটি সংযোগের মাধ্যমে Teacher ও Subject-এর M:N সম্পর্ক বাস্তবায়িত হয়।
Mitu-এর এখন কোনো Teaching রেকর্ড নেই; সম্পর্কের ধরন কেবল বর্তমান
রেকর্ড গুনে নয়, কোন সংযোগ অনুমোদিত তার নিয়ম থেকে নির্ধারিত হয়।


\subsubsection{Student--Course-এর আরেকটি M:N উদাহরণ}

একজন শিক্ষার্থী একাধিক কোর্স নিতে পারে। আবার একই কোর্সে একাধিক শিক্ষার্থী ভর্তি হতে পারে। নিচের Student, Course ও Enrolment টেবিলে এই সংযোগ দেখা যায়।

% Same records used by SQL example 10.2; the junction sits between its parents.
\begin{flushleft}
\begin{tikzpicture}[every node/.style={font=\small,inner sep=0pt,align=left}]
\node (student) {\begin{tabular}{|l|l|}\hline
\multicolumn{2}{|l|}{\textbf{Student}}\\\hline
\chSixHeader StudentID & Name\\\hline
101 & Rupa\\\hline 102 & Nabil\\\hline
\end{tabular}};
\node[anchor=north west] (enrolment) at ([xshift=12mm]student.north east)
{\begin{tabular}{|l|l|}\hline
\multicolumn{2}{|l|}{\textbf{Enrolment}}\\\hline
\chSixHeader StudentID & CourseID\\\hline
101 & C01\\\hline 101 & C02\\\hline 102 & C01\\\hline
\end{tabular}};
\node[anchor=north west] (course) at ([xshift=12mm]enrolment.north east)
{\begin{tabular}{|l|l|}\hline
\multicolumn{2}{|l|}{\textbf{Course}}\\\hline
\chSixHeader CourseID & CourseName\\\hline
C01 & ICT\\\hline C02 & English\\\hline
\end{tabular}};
\draw[bookarrow] (student.east)--node[pos=.2,above=2pt]{1}node[pos=.8,above=2pt]{M}(enrolment.west);
\draw[bookarrow] (course.west)--node[pos=.2,above=2pt]{1}node[pos=.8,above=2pt]{M}(enrolment.east);
\node[anchor=north west,font=\footnotesize] at ([yshift=-2mm]student.south west) {StudentID = PK};
\node[anchor=north west,font=\footnotesize,text width=38mm] at ([yshift=-2mm]enrolment.south west) {StudentID ও CourseID = FK};
\node[anchor=north west,font=\footnotesize] at ([yshift=-2mm]course.south west) {CourseID = PK};
\end{tikzpicture}
\captionof{table}{Student ও Course-এর মাঝখানে Enrolment Junction Table}
\end{flushleft}



Enrolment-এর (101, C01) সারিটি বোঝায় Rupa ICT নিয়েছে। (101, C02) থেকে
বোঝা যায় Rupa English-ও নিয়েছে। আবার C01-এর সঙ্গে 101 ও 102 আছে বলে
ICT-তে Rupa ও Nabil দুজনই রয়েছে। তাই উভয় দিক থেকে একাধিক সংযোগ হয়।


একই Student একই Course-এ একবারই ভর্তি হতে পারলে Enrolment-এর (StudentID, CourseID) জোড়া Composite Primary Key হতে পারে। StudentID একা অনন্য নয়, CourseID-ও একা অনন্য নয়; কিন্তু একই জোড়া আবার বসালে একই ভর্তি দুবার লেখা হবে। একটি ভর্তি বাতিল করলে সেই সংযোগের সারি বাদ যায়; শিক্ষার্থী বা কোর্সের নিজস্ব পরিচয় মুছে ফেলার প্রয়োজন হয় না।

এই নকশায় Student থেকে Enrolment হলো 1:M এবং Course থেকে Enrolment-ও 1:M। অর্থাৎ মূল দুই সত্তার M:N সম্পর্ক মধ্যবর্তী সারণির মাধ্যমে দুটি 1:M সম্পর্কে বাস্তবায়িত হয়। Junction-এর প্রতিটি সারি একজন Student ও একটি Course-এর নির্দিষ্ট সংযোগ বোঝায়; সব Student-এর সঙ্গে সব Course স্বয়ংক্রিয়ভাবে যুক্ত হয় না।

\begin{mistakebox}{এক ফিল্ডে একাধিক ID লিখবে না}

Student টেবিলের একটি CourseID ফিল্ডে \[
C01,C02,C03
\] এভাবে একাধিক মান রাখা ভালো Relational Design নয়। প্রতিটি Student--Course সম্পর্ককে Junction Table-এর আলাদা রেকর্ড
হিসেবে রাখা উচিত।

\end{mistakebox}


% ============================================================
% SQL M:N
% ============================================================

\noindent\begin{minipage}{\linewidth}
\chSixExample{১০.২}{Student, Enrolment ও Course টেবিল থেকে ICT নেওয়া শিক্ষার্থীদের Name দেখাও।}

\noindent
\begin{minipage}[c]{0.60\linewidth}

\begin{verbatim}
SELECT Student.Name
FROM Student, Enrolment, Course
WHERE Student.StudentID =
      Enrolment.StudentID
  AND Enrolment.CourseID =
      Course.CourseID
  AND Course.CourseName = 'ICT'
ORDER BY Student.StudentID;
\end{verbatim}

\end{minipage}
\hfill
\begin{minipage}[c]{0.04\linewidth}
\raggedright
$\Rightarrow$
\end{minipage}
\hfill
\begin{minipage}[c]{0.30\linewidth}
\raggedright

\footnotesize
\textbf{ফল}

\smallskip

\begin{tabular}{|l|}
\hline
\chSixHeader
Name\\
\hline
Rupa\\\hline
Nabil\\
\hline
\end{tabular}

\end{minipage}
\end{minipage}\par\smallskip


\textbf{দুই ধাপে সংযোগ, তারপর বাছাই।} প্রথম শর্তে StudentID দিয়ে
Student থেকে Enrolment-এর সংশ্লিষ্ট সারি মেলে। দ্বিতীয় শর্তে CourseID
দিয়ে ওই Enrolment-এর Course জানা যায়। তাই (101, C01) থেকে Rupa–ICT,
(101, C02) থেকে Rupa–English এবং (102, C01) থেকে Nabil–ICT পাওয়া যায়।
শেষ শর্ত ICT বেছে নেওয়ায় Rupa ও Nabil থাকে; SELECT তাঁদের Name দেখায়।
StudentID আর CourseID ভিন্ন সত্তার পরিচয়—এ দুটিকে সরাসরি সমান না ধরে
Junction Table-এর দুই Foreign Key দিয়ে সংযোগ করতে হয়।


% ============================================================
% ORDER PRODUCT CQ EXAMPLE
% ============================================================

